% F0 -- The whole project on one line, before any vocabulary is introduced.
\begin{figure}[!htbp]
\centering
\fitwidth{%
\begin{tikzpicture}[font=\small,
  st/.style={draw=clrule,fill=clsoft,rounded corners=3pt,align=center,
             inner sep=7pt,text width=3.15cm,minimum height=3.55cm,anchor=north},
  ar/.style={-{Stealth[length=6pt]},clrule,very thick},
  num/.style={circle,draw=clmodel,fill=white,inner sep=1.4pt,
              font=\scriptsize\bfseries,text=clmodel},
]
\node[st] (s1) at (0,0)
  {{\footnotesize\bfseries\color{clmodel} 1. One day of\\ one person}\\[4pt]\footnotesize
   The network writes down which antenna the phone was on, and at what time,
   all day long.};
\node[st] (s2) at (4.35,0)
  {{\footnotesize\bfseries\color{clmodel} 2. Rewritten\\ as a list}\\[4pt]\footnotesize
   That day becomes an ordered list of \emph{antenna, hour} pairs, one line
   per person.};
\node[st] (s3) at (8.70,0)
  {{\footnotesize\bfseries\color{clmodel} 3. A language\\ model reads it}\\[4pt]\footnotesize
   The same kind of model that finishes a sentence, retrained on those lists
   instead of on text.};
\node[st] (s4) at (13.05,0)
  {{\footnotesize\bfseries\color{clmodel} 4. It names the\\ next antenna}\\[4pt]\footnotesize
   Out of the 369 antennas of the study area, it ranks the likeliest one
   first.};

\draw[ar] (s1.east) -- (s2.west);
\draw[ar] (s2.east) -- (s3.west);
\draw[ar] (s3.east) -- (s4.west);
\node[num] at (s1.north west) {1}; \node[num] at (s2.north west) {2};
\node[num] at (s3.north west) {3}; \node[num] at (s4.north west) {4};

% --- the competitor, underneath -----------------------------------------
\node[draw=clbase,dashed,fill=clbase!8,rounded corners=3pt,align=center,
      inner sep=7pt,text width=14.1cm,font=\small,anchor=north]
      (rule) at (6.52,-4.25)
  {\textbf{The thing to beat, on the very same question:} answer
   \emph{``the same antenna as a moment ago''}, every time. No model, no
   training, one line of code.\\[2pt]
   On an ordinary day of records it is right \textbf{69 times out of 100}.};
\draw[ar,clbase] (s4.south) -- (13.05,-4.25);
\end{tikzpicture}}
\caption{What this internship built, in four steps}
\label{fig:overview}
\figtext{%
\figpara{What the picture says.}{Reading left to right: the mobile network
records, all day, which antenna each phone is attached to. I turn one person's
day into a plain ordered list. A small language model, the same kind of program
that predicts the next word of a sentence, is retrained on thousands of those
lists. Asked to continue one of them, it names which antenna the person will be
on next, and it does so by ranking all 369 antennas of the area from the most
to the least likely.}

\figpara{Why a language model, for a problem with no language in it.}{Because
the shape of the problem is the same. A sentence is an ordered list of words
drawn from a fixed vocabulary, and the model's job is to guess the next one. A
day of movement is an ordered list of antennas drawn from a fixed set of 369,
and the job is exactly the same. Nothing else about language is used here.}

\figpara{The dashed box is the real difficulty.}{A competitor with no model at
all, ``the person is still where they were'', is already right about seven times
out of ten, because people mostly stay put. Every number in this report is the
distance between my model and that rule, measured on the same people, on the
same moments of their day.}}
\end{figure}
